\SourceLeaf{334}%
\[
\begin{aligned}
\log t={}&\log\tau+\frac{Mk^4}{12aa(1-ee)}\mathord{.}rr+\tfrac1{24}M\tau\tau\\
&+\frac M{24k^4}(5ee+(4ee-14e^4)ss+5e^4s^4)\beta\beta\\
\log b={}&\log\beta+\frac{Mk^4}{12aa(1-ee)}\mathord{.}rr+\tfrac18M\tau\tau\\
&-\frac M{24k^4}(2+ee-(8ee-14e^4)ss-9e^4s^4)\beta\beta\\
\log l={}&\log\lambda+\tfrac1{24}M\tau\tau-\frac{(1-ee)M}{24k^4}(1-10eess)\beta\beta
\end{aligned}
\]

As is easy to see, all the formulas developed so far assume that the quantities \(t,\tau,b,\beta,l,\lambda\) are expressed in units of the radius. If these quantities are to be expressed in arcseconds and the same symbols are to be retained, the formulas for \(\tau,\beta,\lambda\) (Art. 29) must be given an additional factor \(\frac1\rho\). In equations 9, 10, 11, however, and in the logarithmic equations derived from them, the terms containing \(\tau\tau\) or \(\beta\beta\) must be given an additional factor \(\rho\rho\). Here \(\rho\), as in Arts. 16 and 19 above, denotes the length of an arc subtending one arcsecond in units of the radius. If we also retain \(\mu\) with the meaning used above, namely
\[
\mu=\tfrac1{12}M\rho\rho
\]
and abbreviate as follows
\[
\begin{aligned}
(1)&=\frac{k}{a\rho}\\
(2)&=\frac{k^3}{a(1-ee)\rho}\\
(3)&=\frac{Mk^4}{12aa(1-ee)}\\
(4)&=\frac{\mu}{2k^4}(5ee+(4ee-14e^4)ss+5e^4s^4)\\
(5)&=\frac{\mu}{2k^4}(2+ee-(8ee-14e^4)ss-9e^4s^4)\\
(6)&=\frac{(1-ee)\mu}{2k^4}(1-10eess)\\
(7)&=\tfrac12\mu
\end{aligned}
\]
then the following six formulas give our solution
\[
\begin{aligned}
\tau&=(1)r\sin T\operatorname{tang}B\\
\beta&=(2)r\cos T\\
\lambda&=(1)r\sin T\sec B\\
\log t&=\log\tau+(3)rr+(4)\beta\beta+(7)\tau\tau\\
\log b&=\log\beta+(3)rr-(5)\beta\beta+3(7)\tau\tau\\
\log l&=\log\lambda-(6)\beta\beta+(7)\tau\tau
\end{aligned}
\]
